% =====================================================================
%  IA-SASRec negative-result paper — formatted for the journal
%  "Artificial Intelligence" (ISSN 1561-5359, Oles Honchar Dnipro NU).
%  2-column, 12pt Times New Roman, bilingual UA+EN title/abstract.
%
%  Build:  make pdf
% =====================================================================
\documentclass[twocolumn,12pt,a4paper]{article}

\usepackage[margin=2cm]{geometry}
\usepackage{iftex}
\ifPDFTeX
  % Fallback for pdfLaTeX (Cyrillic via T2A; uses Latin-Modern as Times has no T2A).
  \usepackage[T2A,T1]{fontenc}
  \usepackage[utf8]{inputenc}
  \usepackage[ukrainian,english]{babel}
\else
  % Recommended: xelatex / lualatex with native UTF-8 and a Times clone.
  \usepackage[no-math]{fontspec}
  % "Times New Roman" is a macOS system font and has full Cyrillic
  % coverage; TeX Gyre Termes (the open-source Times clone bundled
  % with TeX Live) does not, so we use the system font on macOS.
  \setmainfont{Times New Roman}
  \usepackage[ukrainian,english]{babel}
\fi

% Forbid algorithmic mid-word hyphenation at line ends. Real hyphens in
% compound names (IA-SASRec-Add, ML-1M, ...) are NOT affected — those
% are controlled by \exhyphenpenalty which we leave at its default.
% Compensate with looser inter-word spacing so paragraphs still fully
% justify to the column width.
\hyphenpenalty=10000
\sloppy
\tolerance=1500
\emergencystretch=3em
\usepackage{amsmath,amssymb}
\usepackage{booktabs}
\usepackage{graphicx}
\usepackage{xcolor}
\usepackage{caption}
\usepackage[style=apa,citestyle=numeric-comp,sortcites=true,sorting=none,backend=biber,natbib=true]{biblatex}
\addbibresource{refs.bib}
\defbibenvironment{bibliography}
  {\enumerate[label=\arabic*.,leftmargin=*,labelsep=0.5em,itemsep=\bibitemsep,parsep=\bibparsep]}
  {\endenumerate}
  {\item}
\usepackage{microtype}
\usepackage{enumitem}
\usepackage{hyperref}
\hypersetup{colorlinks=true, linkcolor=blue!50!black,
            citecolor=blue!50!black, urlcolor=blue!50!black}
\graphicspath{{figures/}}

% Best/second-best result highlighting (used by tables/main_results.tex).
\newcommand{\best}[1]{\textbf{#1}}
\newcommand{\second}[1]{\underline{#1}}

% Caption tuning — body 12pt, captions a notch smaller per sample.
\captionsetup{font=small,labelfont=bf}

% Journal header (top of first page, per SampleArticleEN.pdf).
\usepackage{fancyhdr}
\pagestyle{fancy}
\fancyhf{}
\fancyfoot[R]{\footnotesize \thepage}
\renewcommand{\headrulewidth}{0pt}
\renewcommand{\footrulewidth}{0pt}

\begin{document}

% --- Single-column title block at the top of a two-column document ---
\twocolumn[
\begin{@twocolumnfalse}
\begin{center}
{\footnotesize\bfseries UDC 004.85}\hfill\mbox{}\par
\vspace{0.6em}

% ---- Ukrainian title ----
{\small\itshape Назва статті українською}\par
\vspace{0.2em}
{\bfseries\large ІНТЕНСИВНО-ЗВАЖЕНА САМОУВАГА ДЛЯ ПОСЛІДОВНИХ РЕКОМЕНДАЦІЙ:\par
НЕГАТИВНИЙ РЕЗУЛЬТАТ І АНАЛІЗ МЕХАНІЗМУ}\par
\vspace{0.8em}

% ---- English title ----
{\small\itshape Title of the article in English}\par
\vspace{0.2em}
{\bfseries\large INTENSITY-AWARE SELF-ATTENTION FOR SEQUENTIAL RECOMMENDATION:\par
A NEGATIVE RESULT AND A MECHANISM ANALYSIS}\par
\vspace{1.0em}
\end{center}

% ---- Ukrainian annotation (~1800 chars w/ spaces, per journal spec) ----
% NOTE: native-speaker review recommended before submission. Draft is a
% literal rendering of the English annotation; terminology (увага,
% шар, базова модель, довірчий інтервал, бутстрап) should be checked
% against current Ukrainian recommender-systems literature.
\begin{otherlanguage}{ukrainian}
{\small\noindent
Послідовні рекомендаційні моделі на основі самоуваги відкидають
важливий контрольований сигнал --- \textit{інтенсивність} кожної
взаємодії (оцінку $1$--$5$, кількість годин гри, час перегляду, суму
покупки). Ми перевіряємо гіпотезу, що повернення інтенсивності у
обчислення уваги має покращити топ-$k$ точність пропорційно до
інформативності сигналу. Запропоновано три модифікації SASRec ---
IA-SASRec-Add, IA-SASRec-Mul, IA-SASRec-Val, --- що відповідають
трьом алгебраїчним позиціям, де інтенсивність можна вбудувати у вираз
уваги. Кожен варіант додає один навчальний скаляр $\lambda$ на шар
уваги, так що при $\lambda{=}0$ модель точно збігається з базовою
SASRec. На сімох наборах даних (MovieLens$\times2$, Amazon$\times2$,
Steam$\times3$) і п'ятьох початкових станах на конфігурацію, із
потрійним протоколом значущості (попарний $t$-критерій по
\emph{seed}-ах, попарний бутстрап по користувачах при $B{=}2{,}000$,
мажоритарна згода знаків довірчих інтервалів по \emph{seed}-ах),
жоден варіант не дає робастного покращення. Картина протилежна
гіпотезі: на Steam, де сигнал інтенсивності найбагатший,
спостерігаються найбільші статистично значущі погіршення (до
$-9{,}5\%$ NDCG@10, $p<0{,}05$). Аналіз вивченого $\lambda$ пояснює
провал: $\lambda$ вільно зменшується там, де інтенсивність
неінформативна (до $0{,}07$ на ML-1M), але залишається у діапазоні
$0{,}55$--$0{,}99$ на Steam. Архітектурний ``вихід при $\lambda{=}0$''
доступний у принципі, але не реалізується на практиці. Робимо
висновок, що один скаляр на шар --- надто грубий контроль для
інжекції інтенсивності, та окреслюємо три конкретні архітектурні
альтернативи.

\noindent\textbf{Ключові слова:} послідовні рекомендації;
самоувага; SASRec; увага з урахуванням інтенсивності;
відтворюваність; перевірка значущості; негативний результат;
довірчий інтервал методом бутстрап; аналіз механізму.}
\end{otherlanguage}

\bigskip

% ---- English annotation (from sections/00_abstract.tex) ----
% Section: Abstract
% =====================================================================
% DESCRIPTION
% -----------
% Target length: ~150 words. Must hit five beats in order:
%   1. Hypothesis: per-interaction intensity (rating / playtime / dwell)
%      is richer than binary presence and should help attention-based
%      sequential recommenders.
%   2. What we tried: three drop-in IA-SASRec variants (Add / Mul / Val)
%      that inject a learned per-layer scalar lambda gating the
%      intensity signal into self-attention.
%   3. Protocol: 7 datasets (ML-100K, ML-1M, Amazon Digital Music,
%      Amazon Office Products, Steam-3k/8k/15k), 5 seeds each (4 for
%      ML-1M), Optuna HPO per (dataset, model), evaluated with paired
%      t-test + per-user paired bootstrap + per-seed CI sign agreement.
%   4. Headline: no robust improvement on NDCG@10 / Recall@10 anywhere;
%      Steam — the dataset with the strongest, most continuous intensity
%      signal — shows the LARGEST significant degradations (up to -9.5%
%      NDCG@10, p<0.05).
%   5. Mechanism: learned lambdas remain large (0.5-1.0) precisely where
%      the model is being hurt; the per-layer scalar gate cannot opt out
%      of harmful signal. We argue this rules out simple gating as a
%      sufficient escape hatch and outline what intensity-injection
%      designs should look like instead.
% Lift-from: ia_sasrec.md §1; numeric results re-confirmed against
%            results/eval/*.json.
% =====================================================================
\begin{abstract}
TODO: 150-word abstract following the five-beat structure above.
\end{abstract}


\vspace{1em}
\end{@twocolumnfalse}
]

% --- Body sections (two columns from here) ---
% Section: Introduction
% =====================================================================
% DESCRIPTION
% -----------
% Goal: motivate the intensity-injection hypothesis, state the gap,
% declare a negative result honestly, and preview the Steam paradox so
% the reader knows the paper has a mechanism story (not just a null).
%
% Paragraph plan (~5 paragraphs):
%   P1. Sequential recommenders treat history as a token stream;
%       attention models (SASRec, BERT4Rec) dominate. But every
%       interaction is collapsed to a binary "happened" — the
%       underlying *intensity* (rating, hours played, dwell time) is
%       discarded. That is a lot of supervised signal thrown away.
%   P2. A natural hypothesis: re-introducing per-interaction intensity
%       into self-attention should help, in proportion to how rich the
%       intensity signal is. Steam playtime (continuous, heavy-tailed)
%       should help more than 5-star ratings (discrete, low entropy).
%   P3. We test this with three drop-in modifications to SASRec
%       (Add / Mul / Val), each adding ONE learnable scalar per layer
%       so that lambda=0 collapses exactly to SASRec — the model can
%       opt out if intensity is uninformative.
%   P4. Headline finding: the hypothesis fails. Across 7 datasets and
%       5 seeds per (dataset, model), no variant achieves a robust
%       NDCG@10 / Recall@10 improvement, and the failure is most
%       pronounced on Steam — the OPPOSITE of what the hypothesis
%       predicts. We confirm this with paired t-tests, per-user paired
%       bootstrap CIs, and per-seed sign-agreement.
%   P5. Mechanism: learned lambda stays high (0.55-0.99 on Steam) even
%       as metrics degrade, so the "opt-out" gate fails to fire. The
%       paper's contribution is the documented failure together with
%       the lambda-vs-Delta calibration evidence that simple
%       per-layer scalar gating is too coarse for heavy-tailed
%       intensities.
%
% Reusable prose: ia_sasrec.md §1 (motivation paragraphs) — lift but
% retune the tone from "we propose" to "we hypothesised".
%
% Contributions bullet list at end of section:
%   - Three IA-SASRec variants with explicit lambda fallback.
%   - A 7-dataset, 5-seed evaluation with a dual-criterion significance
%     protocol (paired t + per-user bootstrap + per-seed sign agreement).
%   - A negative result with mechanism analysis via the learned lambda.
%   - Released code, checkpoints, and per-user metric parquets for full
%     reproducibility.
% =====================================================================
\section{Introduction}
\label{sec:intro}

Sequential recommenders treat a user's interaction history as a token
stream, and self-attention models such as SASRec~\citep{kang2018sasrec}
and BERT4Rec~\citep{sun2019bert4rec} dominate the leaderboards. In
the standard formulation, every interaction is collapsed to a binary
``happened'': item $i_j$ was in the user's history at position $j$,
and that is the totality of the signal that reaches the attention
layer. The underlying intensity of each interaction (a
$5$-star rating on MovieLens, hours played on Steam, dwell time or
purchase value in industrial logs) is discarded. From a Shannon
standpoint that is a lot of supervised information thrown away, and
the discard is more notable the richer the intensity signal is:
a binary signal carries one bit; a continuous heavy-tailed playtime
distribution can carry several.

A natural hypothesis follows: re-introducing per-interaction intensity
into the attention computation should help, and should help in
proportion to how informative the underlying signal is. Steam
playtime (continuous, heavy-tailed, several bits per interaction)
should help more than discrete $1$--$5$ ratings on
MovieLens, whose entropy is bounded above by $\log_2 5 \approx 2.3$
bits and is much lower in practice given the well-known rating bias
toward $4$ and $5$.

We test this hypothesis with three drop-in modifications of SASRec
(IA-SASRec-Add, IA-SASRec-Mul, and IA-SASRec-Val), corresponding
to the three algebraic positions in the attention expression where
intensity can be threaded in (pre-softmax additive bias, pre-softmax
multiplicative scale, post-softmax attention reweighting).
Each variant adds exactly one learnable scalar $\lambda$ per
attention layer, initialised to $1.0$, designed so that $\lambda{=}0$
collapses the variant bit-for-bit to vanilla SASRec. The model is
therefore a strict super-set of the baseline: if intensity is
not informative, a fair learner can recover SASRec exactly.

Across seven datasets, spanning the discrete-rating and
continuous-playtime regimes, and five random seeds per
(dataset, model) cell, the hypothesis fails. Not one of the $63$
tested cells ($3$ variants $\times$ $3$ metrics at $k{=}10$ across
$7$ datasets) shows a robust improvement under our dual-criterion
significance protocol: a seed-level paired $t$ and a per-user
paired bootstrap that must both reject, guarded by a per-seed
CI sign-majority on the bootstrap. The single
positive paired-$t$ cell (IA-SASRec-Add on Amazon-Office MRR@10,
$+3.6\%, p{=}0.025$) is suppressed by the per-user bootstrap and
the per-seed sign agreement. Worse, the degradation pattern is the
opposite of the intensity-richness prediction. Steam, the
dataset with the richest signal, hosts the largest and cleanest
significant losses ($-9.5\%$ to $-7.3\%$ NDCG@10/Recall@10/MRR@10
across Steam-3k Mul and Steam-15k Mul/Val, all robust).

The architecture was designed with an explicit fallback
($\lambda{=}0$ recovers SASRec), so the failure has a sharp
mechanism question attached, which we answer through an analysis of
the learned per-layer $\lambda$. The learned $\lambda$ does adapt: it
drops to $0.07$ on ML-1M Val (and the model recovers SASRec), so the
optimisation path is open. But on Steam it stays at $0.55$--$0.99$
across all three variants and all three subset sizes, even as test
metrics degrade. The optimiser can disable intensity and chooses not to. A
single per-layer scalar is not a sufficient fallback for
intensity injection: it cannot down-weight specific intensity values,
and it cannot let some attention heads use intensity while others
ignore it. It also confounds the signal itself with the way we encoded
it: a flat $\lambda$ could mean intensity carries no information, or
that our normalisation choice flattened the information that was there.

\paragraph{Contributions.}
\begin{itemize}
\item Three IA-SASRec variants (Add / Mul / Val) realised as a
unified framework with one learnable scalar $\lambda$ per layer, so
that $\lambda{=}0$ collapses each variant exactly to SASRec.
\item A 7-dataset $\times$ 5-seed evaluation under a dual-criterion
significance protocol (paired-$t$ at the seed level and a per-user
paired bootstrap with $B{=}2{,}000$, guarded by a per-seed CI
sign-majority on the bootstrap) explicitly designed against
the failure modes flagged by prior reproducibility critiques~\citep{ferrari2019,sun2020daisy}.
\item A negative result with mechanism evidence: learned $\lambda$
stays large precisely where the variant hurts the most, and the
``off-switch at $\lambda{=}0$'' fallback is not realised in
practice.
\end{itemize}

% Section: Related Work
% =====================================================================
% DESCRIPTION
% -----------
% Three buckets, each ~one paragraph. Goal is to (a) place IA-SASRec in
% the literature and (b) make the gap concrete: no prior work isolates
% PER-INTERACTION INTENSITY as the injected attention signal.
%
% Bucket A — Sequential attention recommenders:
%   - SASRec \citep{kang2018sasrec}: causal self-attention over the
%     interaction sequence; the baseline.
%   - BERT4Rec \citep{sun2019bert4rec}: bidirectional masked-LM
%     pretraining.
%   - TiSASRec \citep{li2020tisasrec}: injects RELATIVE TIME GAPS
%     between interactions into attention. Closest prior art — but the
%     injected signal is temporal, not interaction strength.
%
% Bucket B — Side-information injection in sequence models:
%   - FDSA \citep{zhang2019fdsa}, S3-Rec \citep{zhou2020s3rec}: inject
%     item-side features (category, attribute) via separate streams or
%     pretraining objectives.
%   - Frame the distinction: these inject *item* metadata; IA-SASRec
%     injects *interaction-level* intensity, which varies across users
%     for the same item.
%
% Bucket C — Reproducibility / benchmarking critiques:
%   - Ferrari Dacrema et al. \citep{ferrari2019} — many neural rec
%     models fail to beat tuned baselines once protocol is fixed.
%   - Sun et al. \citep{sun2020daisy} — single-seed reporting inflates
%     gains. Our triple-gate significance protocol is designed against
%     exactly these failure modes; the Amazon-Office-Products positive
%     cell (Add, MRR@10, +3.6%, p=0.025) would have been a publishable
%     "win" under a one-seed protocol — we show it does not generalise.
%
% refs.bib entries to add (skeleton placeholders in refs.bib):
%   kang2018sasrec, sun2019bert4rec, li2020tisasrec, zhang2019fdsa,
%   zhou2020s3rec, ferrari2019, sun2020daisy, zhao2021recbole,
%   akiba2019optuna.
% =====================================================================
\section{Related Work}
TODO: related-work prose across the three buckets above.

% Section: Method — IA-SASRec
% =====================================================================
% Lifted from ia_sasrec.md §2--4. Three variants, learnable lambda per
% layer (init 1.0), normalisation as a hyperparameter.
% =====================================================================
\section{IA-SASRec: Intensity-Aware Self-Attention}
\label{sec:method}

\subsection{Background: SASRec attention}
\label{sec:method:background}

SASRec~\citep{kang2018sasrec} encodes a user's left-padded history
$s = [i_1, \dots, i_T]$ (sequence length $T$) as embeddings
$E \in \mathbb{R}^{T \times d}$, where $d$ is the embedding dimension.
Each transformer block projects $E$ to queries, keys, and values
$Q, K, V \in \mathbb{R}^{T \times d}$, and computes scaled dot-product
self-attention with a causal mask $M \in \{0, -\infty\}^{T \times T}$:
\begin{equation}
\mathrm{Attention}(Q, K, V) = \mathrm{softmax}\!\left(\tfrac{QK^{\top}}{\sqrt{d}} + M\right) V.
\label{eq:sasrec}
\end{equation}
Only an item's identity (via the embedding lookup) and its position (via the
positional embedding) enter equation~\eqref{eq:sasrec}; the
strength of each interaction (a 5-star rating, hours played,
or dwell time) is not used. IA-SASRec re-introduces that signal.

Let $w \in \mathbb{R}^T$ be the per-position intensity vector for the
user's history (after the normalisation step described below). Define
the intensity matrix $M_W \in \mathbb{R}^{T \times T}$ with
$M_W[q, k] = w_k$, broadcasting intensity along the key axis,
which is the axis the softmax normalises.

\subsection{Three injection points}
\label{sec:method:variants}

The attention expression~\eqref{eq:sasrec} has three algebraic positions
where $M_W$ can be threaded in: as a pre-softmax additive bias, as a
pre-softmax multiplicative scale, or as a post-softmax attention reweighting.
We instantiate all three. Each adds one learnable scalar
$\lambda \in \mathbb{R}$ per attention layer, initialised to $1.0$, so
that $\lambda = 0$ collapses the variant exactly to
SASRec~\eqref{eq:sasrec}. This learnable gate is the central design
discipline of IA-SASRec.

\paragraph{IA-SASRec-Add (additive logit bias).}
Intensity enters as an explicit bias on the raw attention logits:
\begin{equation}
\mathrm{Att}_{\mathrm{Add}} = \mathrm{softmax}\!\left(\tfrac{QK^{\top}}{\sqrt{d}} + \lambda\, M_W + M\right) V.
\label{eq:add}
\end{equation}
A high-intensity key $k$ inflates column $k$ of the logit matrix and is
forced to receive large attention probability, irrespective of semantic
similarity to the query. As $\lambda \to \infty$ the variant approaches
a hard intensity-argmax.

\paragraph{IA-SASRec-Mul (multiplicative logit scaling).}
Instead of an additive bias, the semantic similarity itself is scaled
by intensity:
\begin{equation}
\begin{aligned}
\mathrm{Att}_{\mathrm{Mul}} = \mathrm{softmax}\Big(
  & \tfrac{QK^{\top}}{\sqrt{d}} \odot \big(1 + \lambda(M_W - 1)\big) \\
  & {} + M \Big)\, V,
\end{aligned}
\label{eq:mul}
\end{equation}
with $\odot$ element-wise. The $(w_k - 1)$ parameterisation interpolates
between two endpoints: at $\lambda = 0$ the multiplier is $1$ and the
variant collapses to SASRec; at $\lambda = 1$ the multiplier is exactly
$w_k$, recovering the original hard-gatekeeper formulation. The causal
mask is added after scaling so padded keys remain at $-\infty$.

\paragraph{IA-SASRec-Val (post-softmax attention reweighting).}
The standard softmax computes the attention distribution; intensity then
reweights each key position's attention probability before the value
mix. Equivalently, since the factor depends only on $k$, this is
algebraically identical to pre-scaling each value vector $V_k$ by its
intensity: $\big(A \odot (1{+}\lambda(M_W{-}1))\big)\,V = A\,(D_\lambda V)$,
where $D_\lambda = \mathrm{diag}(1{+}\lambda(w{-}1))$. Hence the name
\emph{Val}: the operation acts at the value-mixing stage of attention.
\begin{equation}
\begin{aligned}
\mathrm{Att}_{\mathrm{Val}} = \bigl[\,
  & \mathrm{softmax}\!\left(\tfrac{QK^{\top}}{\sqrt{d}} + M\right) \\
  & \odot \bigl(1 + \lambda(M_W - 1)\bigr) \bigr]\, V.
\end{aligned}
\label{eq:val}
\end{equation}
Unlike Mul, which acts on raw logits and can be diluted by the
softmax normalisation, Val acts directly on the attention
probabilities that mix the values. We deliberately do not renormalise
the reweighted distribution: any magnitude inflation introduced by
$1{+}\lambda(w_k{-}1) > 1$ is absorbed by the post-attention layer
normalisation~\citep{ba2016layernorm} in the residual block, while the
intensity-driven \emph{direction} change in the output vector --- the
part that actually re-ranks candidates --- passes through unaffected.

Equation~\eqref{eq:val} is a breaking change from a previous formulation
of Val, which multiplied the post-attention context by
query-position intensity. At prediction time only the last query is
read out, so every candidate score was rescaled by the same constant
and the ranking was identical to vanilla SASRec: the variant could
not reorder candidates and is not a meaningful test of intensity
injection.

\subsection{The $\lambda$ fallback}
\label{sec:method:escape}

At $\lambda = 0$ each variant equals SASRec bit-for-bit (verified by
unit tests on the attention math). A fair learner that finds intensity
not informative should therefore drive $\lambda \to 0$ and recover the
baseline. This is the load-bearing prediction we test empirically: if
the fallback holds, learned $\lambda$
should be small wherever IA-SASRec fails to improve over SASRec, and
large only where it helps.

\subsection{Intensity normalisation}
\label{sec:method:norm}

Raw intensities are heterogeneous: Steam hours range from $0.1$ to
$1000+$, while ratings live in $\{1,\dots,5\}$. Feeding raw values
through a softmax causes severe value collapse, so normalisation is
essential rather than cosmetic. We expose four modes as a single
hyperparameter \texttt{intensity\_norm} selected by Optuna~\citep{akiba2019optuna}:
\texttt{log1p\_minmax} (default; per-user $\log(1+w)$ then divide by
the masked max), \texttt{minmax} (per-user with a floor $f = 0.1$ so
the weakest real signal stays distinguishable from padding),
\texttt{zscore}, and \texttt{none} (sanity-check ablation). The floor
$f$ is non-cosmetic: without it the least-intense real item maps to
exactly $0$ and becomes indistinguishable from padding, silently
discarding one interaction per sequence in Mul and Val. After every
mode, padding positions are forcibly zeroed.

\subsection{Parameter cost}
\label{sec:method:cost}

Each variant adds exactly $n_{\mathrm{layers}}$ learnable scalars
total ($2$--$3$ in our configurations), strictly less than any
prior side-information injection. Table~\ref{tab:variants} summarises
the three formulations.

\begin{table*}[t]
\centering
\small
\caption{The three IA-SASRec variants. $\lambda$ is a learnable scalar per layer (init $1.0$); $\lambda = 0$ collapses each variant to SASRec.}
\label{tab:variants}
\begin{tabular}{ll}
\toprule
Variant & Mechanism \\
\midrule
Add & $\mathrm{softmax}\!\left(\dfrac{QK^{\top}}{\sqrt{d}} + \lambda\, M_W + M\right)\! V$ \\[1.2em]
Mul & $\mathrm{softmax}\!\left(\dfrac{QK^{\top}}{\sqrt{d}} \odot \bigl(1 + \lambda(M_W - 1)\bigr) + M\right)\! V$ \\[1.2em]
Val & $\left[\,\mathrm{softmax}\!\left(\dfrac{QK^{\top}}{\sqrt{d}} + M\right) \odot \bigl(1 + \lambda(M_W - 1)\bigr)\right] V$ \\[0.6em]
\bottomrule
\end{tabular}
\end{table*}

% Section: Experimental Setup
% =====================================================================
% DESCRIPTION
% -----------
% Three subsections. The point is to make the negative result
% UNATTACKABLE on protocol grounds — every "you under-tuned the
% variant" or "you got unlucky with seeds" objection must be answered
% here.
%
%   4.1 Datasets (Table 2).
%       Seven datasets chosen to span intensity-signal richness:
%         Discrete low-entropy:  ML-100K-iar, ML-1M, Amazon Digital
%                                Music, Amazon Office Products
%                                (1-5 integer ratings).
%         Continuous heavy-tail: Steam-3k / Steam-8k / Steam-15k
%                                (hours played; subset sizes refer to
%                                top-N most-active users).
%       Columns: |U|, |I|, |interactions|, density, intensity source,
%       leave-one-out test target. Pull numbers from src/data.py + the
%       loaded .inter files; report after the project's standard 5-core
%       filtering.
%
%   4.2 Training & HPO.
%       - Identical Optuna search space per (dataset, model) — see
%         src/hpo.py and src/models/__init__.py for the registry.
%       - HPO objective: validation NDCG@10. Budget: N_TRIALS Optuna
%         trials per (dataset, model). Report the actual trial counts
%         and wall-clock per cell to demonstrate budget parity (this
%         is what reviewers attack first).
%       - Final fit: train with the seed-fixed best_params for up to
%         MAX_EPOCHS with early stopping on validation NDCG@10.
%       - Five seeds {2020, 2021, 2022, 2023, 2024}.
%
%   4.3 Evaluation protocol (the methodological backbone).
%       Triple-gate significance, all three must hold for a result to
%       count as "robust":
%         (i)   Seed-level paired t-test (paired_significance in
%               src/evaluation.py), df = n_seeds - 1, alpha = 0.05.
%               Measures TRAINING-STOCHASTICITY variance.
%         (ii)  Per-user paired bootstrap CI (B=2000 resamples) on
%               aligned per-user metric diffs. Measures
%               TEST-POPULATION variance.
%         (iii) Per-seed CI sign-agreement majority — guards against
%               a single lucky seed driving the pooled bootstrap.
%       Wilcoxon signed-rank is reported for completeness but flagged
%       as INFORMATIONAL ONLY at n=5 (exact p-floor is 0.0625 — cannot
%       reject at alpha=0.05).
%       Cohen's d on per-user diffs reported as the
%       PRACTICAL-significance reality check; with thousands of test
%       users, statistical significance can be trivial.
%
%   4.4 Implementation & reproducibility.
%       Built on RecBole \citep{zhao2021recbole}; HPO via Optuna
%       \citep{akiba2019optuna}. Code, checkpoints, eval JSONs, and
%       per-user metric parquets released; the analysis is reproducible
%       end-to-end from \texttt{notebooks/1\_significance.ipynb}.
% =====================================================================
\section{Experimental Setup}
TODO: 4.1 datasets table, 4.2 training/HPO budget, 4.3 triple-gate
significance protocol, 4.4 implementation note.

% Section: Results
% =====================================================================
% DESCRIPTION
% -----------
% Single primary table + one figure + three narrative paragraphs (one
% per dataset cluster). Do NOT bury the negative cells in an appendix —
% the pattern (worst where signal is richest) IS the story.
%
%   Main table (auto-generate via tables/main_results.tex from
%   notebooks/1_significance.ipynb cell 11 [combined] / cell 12 [robust]):
%     Rows  : 7 datasets (grouped: ML / Amazon / Steam).
%     Cols  : {Add, Mul, Val} x {NDCG@10, Recall@10, MRR@10}.
%     Cells : rel%Δ vs SASRec, paired-t stars (*/**/***), per-seed
%             sign indicator "+k/-k/n", and an asterisk if the pooled
%             bootstrap CI excludes 0.
%     Use \best / \second macros (lifted from
%     docs/protocol_neural_networks.tex) to highlight the (rare)
%     significant cells in both directions — green for positive, red
%     for negative.
%
%   Figure 1 (per-user bootstrap forest plot): one row per
%   (dataset, variant) for NDCG@10. Centre = mean per-user diff,
%   whiskers = 95% bootstrap CI. The visual punchline: most CIs cross
%   zero or sit on the wrong side of it. Render this in Phase 2 of
%   writing from the per_user/*.parquet files.
%
% Narrative paragraphs (~one each):
%   P1 — MovieLens (ML-100K-iar, ML-1M):
%        Every cell n.s. at alpha=0.05 across NDCG@10 / Recall@10;
%        ml-100k shows ~5-8% relative drops (noisy at this scale, n=5).
%        Conclusion: discrete low-entropy rating intensity does
%        nothing — neither helping nor robustly hurting.
%
%   P2 — Amazon (Digital Music, Office Products):
%        Mixed and small. The ONE positive-significant cell in the
%        entire campaign: IA-SASRec-Add on Amazon-Office-Products
%        MRR@10 (+3.6%, p=0.025). NDCG@10 also +2.5% but only p=0.083.
%        Mul on Digital Music is significantly WORSE on NDCG@10
%        (-2.7%, p=0.012) and Recall@10 (-2.9%, p=0.021). Val on
%        Office Products is significantly worse (Recall@10 -6.1%,
%        p=0.005). Treat the one positive as a single-cell
%        observation, not a finding — it is exactly the kind of result
%        a triple-gate protocol is designed to NOT over-claim.
%
%   P3 — Steam (3k / 8k / 15k):
%        The "hypothesis-killer" cluster. Every variant negative on
%        NDCG@10 across all three subset sizes. Steam-3k Mul -9.5%
%        (p=0.027). Steam-15k Mul -5.0% (p=0.009), Val -6.2%
%        (p=0.005), Recall@10 -4.5% / -5.1% with p<0.01. The
%        degradation is the OPPOSITE of the intensity-richness
%        prediction and is the cleanest signal of method failure in
%        the campaign.
%
% Cross-reference into §6 for the mechanism analysis.
% =====================================================================
\section{Results}
TODO: insert main results table (tables/main_results.tex), forest plot
(figures/forest_ndcg10.pdf), and three narrative paragraphs (MovieLens,
Amazon, Steam) per the description above.

% Section: Analysis — Why does intensity hurt?
% =====================================================================
% DESCRIPTION
% -----------
% The mechanism story. Without this section the paper is just "we
% tried, it didn't work". With this section it becomes
% "we tried, it didn't work, and here's the architectural reason —
% which the field should know before building variant #47."
%
%   6.1 lambda calibration figure (Figure 2).
%       Source plot: results/lambda_vs_delta.png — copy into
%       paper/figures/. Each point = one (dataset, variant) pair.
%       X = mean learned lambda across seeds. Y = mean test-NDCG@10 Δ
%       vs SASRec across seeds. Reading:
%         - Top-right quadrant: model uses intensity AND gains.
%           Predicted under the hypothesis; observed: empty / sparse.
%         - Bottom-right quadrant: model uses intensity AND loses.
%           Observed: Steam clusters here with lambda 0.55-0.99.
%         - Origin: model opts out (lambda~0) AND recovers SASRec.
%           Observed: a few ML-1M points sit near origin (Val
%           lambda=0.07).
%       The plot's punchline: there is no positive correlation
%       between lambda and Δ. The escape-hatch story FAILS in practice
%       — high lambda does not predict gain (and on Steam, predicts
%       loss).
%
%   6.2 The Steam paradox.
%       Steam has the strongest intensity signal yet the worst
%       outcomes. Two non-mutually-exclusive hypotheses worth probing:
%         (a) Distribution shape: playtime is heavy-tailed
%             (log-normal-ish). After log1p_minmax, a 100h game vs.
%             a 2h game collapses to a narrow band — the model is
%             forced to attend to a near-constant signal that adds
%             variance without information. To diagnose: plot
%             histogram of post-normalisation intensity per dataset
%             (data is in the .inter files). If Steam's distribution
%             is near-constant after norm, the diagnosis lands.
%         (b) Popularity confound: high-playtime games are also
%             popular; item-embedding learning already absorbs that
%             signal, so the extra channel only adds variance.
%             To diagnose: per-dataset Spearman correlation between
%             per-interaction intensity and item popularity rank.
%
%   6.3 Ablation: lambda = 0 vs. lambda = 1 vs. learned.
%       TODO (FLAG IN PAPER): this ablation has NOT been run yet. It
%       is the single most important pre-submission task because it
%       directly tests the "escape-hatch" claim. Expected outcome on
%       Steam: lambda=0 matches SASRec; lambda=1 matches IA-SASRec;
%       learned sits between — confirming that the model could opt
%       out but doesn't, because there is no usable gradient signal
%       pointing it back to lambda=0.
%
%   6.4 Implication for IA-SASRec design.
%       A SINGLE scalar per layer is too coarse:
%       - It cannot down-weight specific intensity values (e.g. mute
%         the 100h tail while keeping the moderate range).
%       - It cannot let some heads use intensity while others ignore
%         it (lambda is shared across heads in our implementation).
%       - It conflates "intensity is informative" with "intensity, in
%         its current normalised form, is informative" — confounding
%         the architectural claim with the preprocessing choice.
% =====================================================================
\section{Analysis}
TODO: 6.1 lambda-vs-Delta figure + reading, 6.2 Steam paradox probes,
6.3 lambda ablation (mark as REQUIRED before submission), 6.4
architectural implications.

% Section: Discussion
% =====================================================================
% DESCRIPTION
% -----------
% Pull the negative result up to the level of general lessons for the
% sequential-recommender community. Don't over-claim — every sentence
% should be defensible from the data in §5-6.
%
% Three threads:
%
%   7.1 What the result says about intensity as a signal.
%       Intensity is NOT free supervision. The implicit assumption in
%       a lot of "side information helps" literature is that any
%       additional channel can only add information. Our results
%       falsify that assumption for per-interaction intensity under
%       scalar gating: the channel can carry NEGATIVE information
%       once optimisation locks lambda > 0 without a gradient path
%       back to 0.
%
%   7.2 What the result says about benchmarking practice.
%       Under a single-seed protocol with no per-user CIs, the
%       Amazon-Office-Products IA-SASRec-Add MRR@10 +3.6% (p=0.025)
%       would be the headline. Our dual-criterion protocol surfaces it
%       as a single-cell observation: it does not generalise across
%       datasets or even across metrics on the same dataset (the
%       same configuration is n.s. on NDCG@10 and Recall@10). This
%       is exactly the failure mode flagged by Ferrari Dacrema et al.
%       and Sun et al. — and a working example of the protocol
%       catching it.
%
%   7.3 Threats to validity.
%       - Optuna budget: were variants under-tuned? Report the trial
%         counts and final-epoch counts side-by-side; the SASRec
%         best_params are not systematically deeper / wider than the
%         IA-SASRec ones, suggesting budget parity is real.
%       - Single intensity normalisation FAMILY per dataset
%         (best of {none, minmax, log1p_minmax, zscore} per
%         best_params). A richer normalisation search might help on
%         Steam; we acknowledge this and leave it as future work.
%       - The lambda=0 ablation in §6.3 has not been run; without it,
%         the "fallback fails" claim rests on observing high
%         learned lambda + degraded metrics, but cannot exclude the
%         counter-narrative "lambda=0 would also have degraded".
%       - Five-seed paired t is conservative (df=4); we mitigate with
%         the per-user bootstrap, which has thousands of units.
% =====================================================================
\section{Discussion}
\label{sec:discussion}

\subsection{Empirical validation of the decoupled-fusion hypothesis}

Our results directly confirm the structural warnings raised by the
decoupled-fusion literature (NOVA~\citep{liu2021novabert},
DIF-SR~\citep{xie2022difsr}). Because IA-SASRec forces an invasive,
mid-attention fusion through a shared scalar gate, it sits in the
architectural region these works predict will fail. The empirical
verdict matches the prediction: across the $252$ (variant, metric,
dataset, cutoff) cells we test (3 variants $\times$ 3 metrics
$\times$ 7 datasets at $k\in\{10, 20, 50, 100\}$), not a single cell
yields a robust improvement, and the largest, cleanest degradations
appear on Steam, the dataset with the strongest and most continuous
intensity signal, where representation degeneration of the kind
flagged by~\citep{fan2022stosa,qiu2022duorec} is most likely to bite.

\subsection{Intensity is not free supervision}

Much of the side-information literature for sequential recommenders
operates under an implicit additivity assumption: any extra channel
threaded into the model can only add information, since the model
``can ignore it.'' Our results rule out that assumption for
per-interaction intensity under scalar gating. Once the optimiser has
settled into a non-zero $\lambda$, the gradient with respect to
training-set cross-entropy does not point back toward $\lambda{=}0$,
even when held-out generalisation degrades. The channel carries
negative information in the sense that fitting it as instructed
makes the model worse. ``The model can ignore it'' is a claim about
expressivity, not about what the optimiser will actually do.

\subsection{Lesson for benchmarking practice}

Under a single-seed protocol without per-user confidence intervals,
the headline of this paper would have been ``IA-SASRec-Add improves
MRR@10 on Amazon-Office Products by $3.6\%$ ($p{=}0.025$).'' The same
configuration is non-significant on the other two metrics of the same
dataset and on every other dataset, but no single-seed table would
have surfaced that. Our dual-criterion protocol flags the cell as
non-robust automatically; this is exactly the failure mode that prior
reproducibility critiques~\citep{ferrari2019,sun2020daisy} have argued
is common in neural recommender benchmarking. We consider this a
working example of the protocol
doing its job, and we argue that it should be the default rather than
the exception when reporting small-sample sequential-recommender
results.

\subsection{Threats to validity}

\paragraph{HPO budget.}
We use the same $25$-trial Optuna budget for SASRec and for each
IA-SASRec variant, with the same backbone search space and an extra
categorical for the normalisation choice on IA-SASRec. The variants
therefore receive strictly more hyperparameter freedom than
SASRec at no extra cost. The selected configurations are not
systematically deeper or wider for SASRec than for the variants, so
``under-tuned'' is not a plausible explanation of the degradations.

\paragraph{Normalisation family.}
Optuna selects a single normalisation per (dataset, model). A richer
parametric normalisation (e.g.\ learned bin boundaries on the playtime
quantile) might rescue Steam; we leave this to future work and flag
it as the most likely route to ``intensity helps after all'' on
heavy-tailed signals.

\paragraph{No controlled $\lambda{=}0$ run.}
We did not pin $\lambda{=}0$ as a control, so the falsification of the
fallback in our analysis rests on observed
dynamics (large learned $\lambda$ coexists with degraded metrics; the
optimiser can drive $\lambda$ small elsewhere). It does not
strictly exclude the counter-narrative ``$\lambda{=}0$ would also
have degraded on Steam''; we phrase the lesson in
the Conclusion accordingly.

\paragraph{Conservatism of paired-$t$ at $n{=}5$.}
With $\mathrm{df}{=}4$ the paired-$t$ has low power. We mitigate by
also requiring the per-user bootstrap (effective $n$ in the thousands)
and the per-seed sign-majority gate. The cost is that some
small-magnitude positive effects may be missed; the benefit is that
the cells we do flag as robust are unlikely to be artifacts of
seed or test-population sampling.

% Section: Conclusion
% =====================================================================
% DESCRIPTION
% -----------
% ~100 words. Three sentences:
%   (1) Recap: we tested three drop-in intensity-injection variants of
%       SASRec across 7 datasets and 5 seeds with a triple-gate
%       significance protocol; none deliver a robust improvement, and
%       the failure is largest on the dataset with the strongest
%       intensity signal.
%   (2) Lesson: a single per-layer scalar lambda is insufficient as an
%       escape hatch — the learned lambda stays large where the model
%       is being hurt, so the architectural opt-out predicted at
%       lambda=0 is not realised in practice.
%   (3) Future directions worth trying instead:
%         - token-level (per-position) intensity gating rather than
%           per-layer scalar;
%         - mixture-of-experts on intensity bins (e.g. discrete buckets
%           by post-norm percentile);
%         - intensity as a side-token in the sequence rather than an
%           attention modifier (cheaper, decouples the gradient).
%   Plus the standard reproducibility pointer: code, checkpoints,
%   per-user metrics are released.
% =====================================================================
\section{Conclusion}
TODO: 100-word recap + lesson + three future-work pointers per the
plan above.


\printbibliography

\end{document}
